\section*{附錄五\quad 索\qquad 引}
\addcontentsline{toc}{section}{附錄五\texorpdfstring{\quad}{ }索引}
\label{sec:b4_a5}

\setlength{\columnsep}{2em}
\begin{multicols}{2}

\begin{center}
\textbf{\large 一\quad 劃}
\end{center}

\noindent 一般座標系： \hfill I-2\par
\noindent\hspace*{1.5em} $S$-座標（一般座標） \hfill I-1\par
\noindent\hspace*{1.5em} $S_0$-座標（自然座標） \hfill I-2\par
\noindent $\Delta$ \hfill IV-1\par
\noindent $\delta$ \hfill IV-1\par
\noindent $n$ 方根：\par
\noindent\hspace*{1.5em} 複數的 $n$ 方根 \hfill V-2 丁\par
\noindent 0 與 1 的特性 \hfill V-1\par

\vspace{0.8em}
\begin{center}
\textbf{\large 二\quad 劃}
\end{center}

\noindent 二次曲線：\par
\noindent\hspace*{1.5em} 特異二次曲線 \hfill IV-1\par
\noindent\hspace*{1.5em} 非特異二次曲線 \hfill IV-1\par
\noindent 二次式：\par
\noindent\hspace*{1.5em} 二次式的定義 \hfill IV-2 甲\par
\noindent\hspace*{1.5em} 二次式的標準化 \hfill IV-4 乙\par
\noindent\hspace*{1.5em} 橢圓二次式 \hfill IV-5 習題\par
\noindent\hspace*{1.5em} 正定橢圓二次式 \hfill IV-5 習題\par
\noindent\hspace*{1.5em} 負定橢圓二次式 \hfill IV-5 習題\par
\noindent $\arg$（幅角） \hfill V-2 丙\par

\vspace{0.8em}
\begin{center}
\textbf{\large 三\quad 劃}
\end{center}

\noindent 三角恆等式 \hfill V-3 甲、乙\par
\noindent 三角恆等式基本定理 \hfill V-3 丙\par
\noindent 基本三角恆等式 \hfill V-3 甲\par
\noindent 三點共線的條件\par
\noindent\hspace*{1.5em} （複數表示） \hfill V-2 丁\par
\noindent 三垂線定理 \hfill II-3\par
\noindent $\text{AU}$（天文單位） \hfill III-1 乙\par
\noindent $H$ \hfill IV-1\par
\noindent $T^-$（負類錐線） \hfill II-4 甲\par
\noindent $T^0$（零類錐線） \hfill II-4 乙\par
\noindent $T^+$（正類錐線） \hfill II-4 丙\par

\vspace{0.8em}
\begin{center}
\textbf{\large 四\quad 劃}
\end{center}

\noindent 不變性 \hfill IV-1\par
\noindent 不變量 \hfill IV-1\par
\noindent 不變量定理 \hfill IV-3\par
\noindent 天文單位（AU） \hfill III-1 乙\par
\noindent 天體運動的軌跡 \hfill III-1 甲\par
\noindent 中心軸（圓錐的） \hfill II-3\par
\noindent 切線：\par
\noindent\hspace*{1.5em} 球的切線 \hfill II-2\par
\noindent\hspace*{1.5em} 圓的切線 \hfill III-5 乙\par
\noindent\hspace*{1.5em} 橢圓的切線 \hfill III-4 甲\par
\noindent\hspace*{1.5em} 雙曲線的切線 \hfill III-4 乙\par
\noindent\hspace*{1.5em} 拋物線的切線 \hfill III-4 丙\par
\noindent 切線長 \hfill II-2\par
\noindent 切點 \hfill II-2，III-4\par
\noindent 五點決定一二次曲線 \hfill IV-6\par
\noindent 分配律（複數系） \hfill V-1 甲\par
\noindent 半角公式 \hfill V-3 乙\par

\vspace{0.8em}
\begin{center}
\textbf{\large 五\quad 劃}
\end{center}

\noindent 正交：\par
\noindent\hspace*{1.5em} 正交基底 \hfill I-4 甲\par
\noindent\hspace*{1.5em} 正交座標變換 \hfill I-5\par
\noindent\hspace*{1.5em} 正交線性映射 \hfill I-4\par
\noindent 兩直線或曲線正交 \hfill III-6 註\par
\noindent 正性基底 \hfill IV-3\par
\noindent 正性標準化 \hfill IV-3\par
\noindent 正類錐線 \hfill II-4\par
\noindent 正定橢圓二次式 \hfill IV-5 習題\par
\noindent 四點決定一拋物線 \hfill IV-6\par
\noindent 四點共線或共圓的條件 \hfill V-2 丁\par
\noindent 平方根 \hfill V-1，V-2\par
\noindent 平面上的線性映射 \hfill I-3\par
\noindent 平移： \hfill I-2 甲\par
\noindent\hspace*{1.5em} 平移向量 \hfill I-2 甲\par
\noindent\hspace*{1.5em} 平移座標變換 \hfill I-2 甲\par
\noindent 母線（圓錐的） \hfill II-3\par
\noindent 主軸定理 \hfill IV-3\par
\noindent 立體投影 \hfill V-4 甲\par
\noindent 可逆性（複數系的） \hfill V-1 甲\par
\noindent 交換律（複數系的） \hfill V-1 乙\par

\vspace{0.8em}
\begin{center}
\textbf{\large 六\quad 劃}
\end{center}

\noindent 自然：\par
\noindent\hspace*{1.5em} 自然座標系（$S_0$-座標系） \hfill I-3\par
\noindent\hspace*{1.5em} 自然矩陣 \hfill I-3 丙\par
\noindent\hspace*{1.5em} 自然特徵方程式 \hfill IV-2 乙\par
\noindent 共軛：\par
\noindent\hspace*{1.5em} 共軛軸（雙曲線的） \hfill III-2 乙\par
\noindent\hspace*{1.5em} 共軛雙曲線 \hfill III-2 乙\par
\noindent\hspace*{1.5em} 共軛複數 \hfill V-2 戊\par
\noindent 共軸圓系： \hfill 附錄\par
\noindent\hspace*{1.5em} 共軛的共軸圓系 \hfill 附錄\par
\noindent 同心圓系 \hfill 附錄\par
\noindent 有心錐線 \hfill IV-5\par
\noindent 光學性質：\par
\noindent\hspace*{1.5em} 橢球面鏡的光學性質 \hfill III-4 甲\par
\noindent\hspace*{1.5em} 雙曲面鏡的光學性質 \hfill III-4 乙\par
\noindent\hspace*{1.5em} 拋物面鏡的光學性質 \hfill III-4 丙\par
\noindent 向量的映像 \hfill I-3 乙\par

\vspace{0.8em}
\begin{center}
\textbf{\large 七\quad 劃}
\end{center}

\noindent 拋物線：\par
\noindent\hspace*{1.5em} 拋物線軸 \hfill III-3 甲\par
\noindent\hspace*{1.5em} 拋物線無漸近線 \hfill III-3 丙\par
\noindent\hspace*{1.5em} 四點決定一拋物線 \hfill IV-6\par
\noindent 拋物線的：\par
\noindent\hspace*{1.5em} 拋物線的定義 \hfill III-3 甲\par
\noindent\hspace*{1.5em} 拋物線的標準式 \hfill II-1 甲，III-3 甲\par
\noindent\hspace*{1.5em} 拋物線的離心率 \hfill II-4 乙，III-3 乙\par
\noindent\hspace*{1.5em} 拋物線的頂點 \hfill III-3 甲\par
\noindent\hspace*{1.5em} 拋物線的準線 \hfill II-5，III-3 甲\par
\noindent\hspace*{1.5em} 拋物線的機械作圖 \hfill III-3 甲\par
\noindent\hspace*{1.5em} 拋物線的焦點 \hfill II-1 甲，III-3 甲\par
\noindent\hspace*{1.5em} 拋物線的焦距 \hfill III-3 甲\par
\noindent\hspace*{1.5em} 拋物線的焦半徑 \hfill III-3 甲\par
\noindent\hspace*{1.5em} 拋物線的焦弦 \hfill III-3 習題\par
\noindent\hspace*{1.5em} 拋物線的正焦弦 \hfill III-3 習題\par
\noindent\hspace*{1.5em} 拋物線的特徵性質 \hfill II-1 甲\par
\noindent\hspace*{1.5em} 拋物線的切線 \hfill III-4 丙\par
\noindent 拋物面： \hfill II-4\par
\noindent\hspace*{1.5em} 拋物面鏡 \hfill II-1 甲，III-4 丙\par
\noindent\hspace*{1.5em} 拋物面鏡的光學性質 \hfill III-4 丙\par
\noindent 里曼（Riemann）球面： \hfill V-4 甲\par

\vspace{0.8em}
\begin{center}
\textbf{\large 八\quad 劃}
\end{center}

\noindent 固有：\par
\noindent\hspace*{1.5em} 固有值 \hfill IV-2 乙\par
\noindent\hspace*{1.5em} 固有向量 \hfill IV-2 乙\par
\noindent 長軸： \hfill III-1 甲\par
\noindent\hspace*{1.5em} 長軸半長 \hfill III-1 甲\par
\noindent 長度壓縮性 \hfill I-3 戊\par
\noindent 空間：\par
\noindent\hspace*{1.5em} 空間壓縮性 \hfill I-3 戊\par
\noindent\hspace*{1.5em} 空間的二次幾何 \hfill IV-4 甲\par
\noindent 定心剛性運動 \hfill I-4 甲\par
\noindent 非特異二次曲線 \hfill IV-1\par
\noindent 和化積公式 \hfill V-3 乙\par
\noindent 函數的線性條件 \hfill I-3 乙\par

\vspace{0.8em}
\begin{center}
\textbf{\large 九\quad 劃}
\end{center}

\noindent 負性基底 \hfill IV-3\par
\noindent 負性標準化 \hfill IV-3\par
\noindent 負類錐線 \hfill II-4\par
\noindent 映像：\par
\noindent\hspace*{1.5em} 向量的映像 \hfill I-3 乙\par
\noindent 相似：\par
\noindent\hspace*{1.5em} 橢圓的相似 \hfill III-1 乙\par
\noindent 相對性原理 \hfill I-3 戊\par
\noindent 保長（函數） \hfill I-4 甲\par
\noindent 保角（函數） \hfill I-4 丁習題\par

\vspace{0.8em}
\begin{center}
\textbf{\large 十\quad 劃}
\end{center}

\noindent 座標變換：\par
\noindent\hspace*{1.5em} 一般性的座標變換 \hfill I-2 乙\par
\noindent\hspace*{1.5em} 正交座標變換 \hfill I-5\par
\noindent\hspace*{1.5em} 平移座標變換 \hfill I-2 甲\par
\noindent\hspace*{1.5em} 剛性座標變換 \hfill I-5\par
\noindent\hspace*{1.5em} 旋轉所對應的座標變換\par
\noindent\hspace*{3.0em} （$D=1$） \hfill I-5\par
\noindent\hspace*{1.5em} 鏡射所對應的座標變換\par
\noindent\hspace*{3.0em} （$D=-1$） \hfill I-5\par
\noindent 矩陣： \hfill I-3 丙\par
\noindent\hspace*{1.5em} 自然矩陣（線性映射相關\par
\noindent\hspace*{1.5em} 於座標系的矩陣表示） \hfill I-3 丙\par
\noindent 剛性運動： \hfill I-4 甲\par
\noindent\hspace*{1.5em} 定心剛性運動 \hfill I-4 甲\par
\noindent 剛性座標變換： \hfill I-5\par
\noindent\hspace*{1.5em} 同原點剛性座標變換 \hfill I-5\par
\noindent 特異二次曲線 \hfill IV-1\par
\noindent 特徵性質：\par
\noindent\hspace*{1.5em} 橢圓的特徵性質 \hfill II-1 乙\par
\noindent\hspace*{1.5em} 雙曲線的特徵性質 \hfill II-1 乙\par
\noindent\hspace*{1.5em} 拋物線的特徵性質 \hfill II-1 甲\par
\noindent 特徵方程式： \hfill IV-2 乙\par
\noindent\hspace*{1.5em} 自然特徵方程式（與\par
\noindent\hspace*{1.5em} 上同義） \hfill IV-2 乙\par
\noindent 根平面 \hfill II-2 習題\par
\noindent 效用函數 \hfill III-4 丙\par
\noindent 馬腰 \hfill IV-4\par
\noindent 馬脊 \hfill IV-4\par
\noindent 高斯（Gauss）平面 \hfill V-2 甲\par
\noindent 倍角公式 \hfill V-3 乙\par

\vspace{0.8em}
\begin{center}
\textbf{\large 十一劃}
\end{center}

\noindent 基底：\par
\noindent\hspace*{1.5em} 基底向量 \hfill I-3 甲\par
\noindent\hspace*{1.5em} 正交基底 \hfill I-4 乙\par
\noindent\hspace*{1.5em} 正性基底 \hfill IV-3\par
\noindent\hspace*{1.5em} 負性基底 \hfill IV-3\par
\noindent 基本三角恆等式 \hfill V-3 甲\par
\noindent 基準點 \hfill I-1\par
\noindent 球： \hfill II-2\par
\noindent\hspace*{1.5em} 球面 \hfill II-2\par
\noindent\hspace*{1.5em} 球面投影 \hfill V-4 甲\par
\noindent\hspace*{1.5em} 球心 \hfill II-2\par
\noindent\hspace*{1.5em} 球的切線 \hfill II-2\par
\noindent\hspace*{1.5em} 球的割線定理 \hfill II-2\par
\noindent\hspace*{1.5em} 球的轉動 \hfill V-4 丙\par
\noindent\hspace*{1.5em} 球面映射與平面映射的\par
\noindent\hspace*{1.5em} 對應 \hfill V-4 乙\par
\noindent\hspace*{1.5em} 球外一點至球的切線等長 \hfill II-2\par
\noindent 頂角（圓錐的） \hfill II-3\par
\noindent 頂點：\par
\noindent\hspace*{1.5em} 圓錐的頂點 \hfill II-3\par
\noindent\hspace*{1.5em} 橢圓的頂點 \hfill III-1 甲\par
\noindent\hspace*{1.5em} 雙曲線的頂點 \hfill III-2 甲\par
\noindent\hspace*{1.5em} 拋物線的頂點 \hfill III-3 甲\par
\noindent 旋轉： \hfill I-4 丙\par

\vspace{0.8em}
\begin{center}
\textbf{\large 十二劃}
\end{center}

\noindent 焦點： \hfill II-1 甲，III-1 甲\par
\noindent\hspace*{1.5em} III-2 甲，III-3 甲\par
\noindent 焦半徑： \hfill III-1 甲，III-2 甲\par
\noindent\hspace*{1.5em} III-3 甲\par
\noindent 焦半徑性質 \hfill II-4 甲，丙\par
\noindent 焦弦：（拋物線的） \hfill III-3 習題\par
\noindent\hspace*{1.5em} 正焦弦 \hfill III-3 習題\par
\noindent 短軸： \hfill III-1 甲\par
\noindent\hspace*{1.5em} 短軸半長 \hfill III-1 甲\par
\noindent 割線： \hfill II-2\par
\noindent\hspace*{1.5em} 球的割線定理 \hfill II-2\par
\noindent 極式 \hfill V-2 丙\par
\noindent 極限軌跡 \hfill II-1 乙\par
\noindent 虛數：（純虛數） \hfill V-1\par
\noindent\hspace*{1.5em} 虛數單位 \hfill V-1\par
\noindent\hspace*{1.5em} 虛部（複數的） \hfill V-1\par
\noindent\hspace*{1.5em} 虛根（imaginary root） \hfill V-2 戊\par
\noindent 虛橢圓 \hfill IV-1\par
\noindent 等溫線 \hfill IV-4\par
\noindent 等軸雙曲線 \hfill III-2 乙\par
\noindent 等價的二組資料 \hfill III-2 乙\par
\noindent 等切軸 \hfill 附錄\par
\noindent 準線 \hfill II-5，III-1 丙\par
\noindent 嵌射（injection） \hfill I-5\par
\noindent 無差異曲線 \hfill III-4 丙\par
\noindent 棣美弗定理 \hfill V-2 丙\par
\noindent 結合律（複數系） \hfill V-1 乙\par
\noindent 幅角（輻角） \hfill V-2 丙\par

\vspace{0.8em}
\begin{center}
\textbf{\large 十三劃}
\end{center}

\noindent 圓：\par
\noindent\hspace*{1.5em} 圓系 \hfill 附錄\par
\noindent\hspace*{1.5em} 圓的方程式 \hfill III-5 甲\par
\noindent\hspace*{1.5em} 圓的標準式 \hfill II-1 乙\par
\noindent\hspace*{1.5em} 兩圓的正交 \hfill III-6\par
\noindent\hspace*{1.5em} 共軸圓系 \hfill 附錄\par
\noindent\hspace*{1.5em} 對於圓的鏡射 \hfill III-6，V-2 戊\par
\noindent 圓錐：\par
\noindent\hspace*{1.5em} 圓錐面 \hfill II-3\par
\noindent\hspace*{1.5em} 圓錐曲線（錐線） \hfill II-3\par
\noindent\hspace*{1.5em} 圓錐的頂點 \hfill II-3\par
\noindent\hspace*{1.5em} 圓錐的頂角 \hfill II-3\par
\noindent\hspace*{1.5em} 圓錐的中心軸 \hfill II-3\par
\noindent\hspace*{1.5em} 圓錐的母線 \hfill II-3\par
\noindent 零類錐線 \hfill II-4\par
\noindent 貫軸（雙曲線的） \hfill III-2 甲\par
\noindent 預算線 \hfill III-4 丙\par
\noindent 跡 \hfill IV-3\par

\vspace{0.8em}
\begin{center}
\textbf{\large 十四劃}
\end{center}

\noindent 對稱：\par
\noindent\hspace*{1.5em} 對稱軸（鏡射軸） \hfill I-4 丙\par
\noindent\hspace*{1.5em} 對稱映射（鏡射） \hfill I-4 丙\par
\noindent\hspace*{1.5em} 對稱線性映射 \hfill IV-2 乙\par
\noindent\hspace*{1.5em} 對稱自然矩陣\par
\noindent\hspace*{1.5em} （前者的矩陣表示） \hfill IV-2 乙\par
\noindent\hspace*{1.5em} 對一點對稱 \hfill I-4 丁\par
\noindent\hspace*{1.5em} 對一直線對稱 \hfill I-4 丁\par
\noindent\hspace*{1.5em} 對原點對稱 \hfill II-1 甲註\par
\noindent\hspace*{1.5em} 對 $Y$ 軸對稱 \hfill II-1 甲註\par
\noindent 對射 \hfill I-5 附註\par
\noindent 複數： \hfill V-1 甲\par
\noindent\hspace*{1.5em} 複數系 \hfill V-1 甲\par
\noindent\hspace*{1.5em} 複數平面 \hfill V-2 甲\par
\noindent\hspace*{1.5em} 複數的極式 \hfill V-2 丙\par
\noindent\hspace*{1.5em} 複數的幅角 \hfill V-2 丙\par
\noindent\hspace*{1.5em} 複數的 $n$ 方根 \hfill V-2 丁\par
\noindent\hspace*{1.5em} 複數的平方根 \hfill V-1，V-2 丁\par
\noindent\hspace*{1.5em} 共軛複數 \hfill V-2 戊\par
\noindent 複角公式 \hfill V-3 甲\par
\noindent 無差異曲線 \hfill III-4 丙\par
\noindent 像集 \hfill I-5 附註\par
\noindent 像原 \hfill I-5 附註\par
\noindent 截角 \hfill II-3\par
\noindent 漸近線：\par
\noindent\hspace*{1.5em} 雙曲線的漸近線 \hfill III-2 乙\par
\noindent\hspace*{1.5em} 拋物線無漸近線 \hfill III-3 丙\par

\vspace{0.8em}
\begin{center}
\textbf{\large 十五劃}
\end{center}

\noindent 線性映射：\par
\noindent\hspace*{1.5em} 線性映射的意義 \hfill I-3 甲\par
\noindent\hspace*{1.5em} 線性映射的矩陣表示 \hfill I-3 丙\par
\noindent\hspace*{1.5em} 線性映射的聯立式 \hfill I-3 丁\par
\noindent\hspace*{1.5em} 線性映射與座標變換\par
\noindent\hspace*{1.5em} 的區別 \hfill I-3 甲注意\par
\noindent\hspace*{1.5em} 線性映射的代數表示 \hfill I-3 丙\par
\noindent\hspace*{1.5em} 平面上的線性映射 \hfill I-3\par
\noindent\hspace*{1.5em} 正交線性映射 \hfill I-4 乙\par
\noindent\hspace*{1.5em} 對稱線性映射 \hfill IV-2 乙\par
\noindent 線性組合： \hfill I-3 乙\par
\noindent 線性條件：\par
\noindent\hspace*{1.5em} 函數的線性條件 \hfill I-3 乙\par
\noindent 標準式：\par
\noindent\hspace*{1.5em} 圓的標準式 \hfill II-1 乙\par
\noindent\hspace*{1.5em} 橢圓的標準式 \hfill III-1 甲\par
\noindent\hspace*{1.5em} 雙曲線的標準式 \hfill III-2 甲\par
\noindent\hspace*{1.5em} 拋物線的標準式 \hfill II-1 甲，III-3 甲\par
\noindent 標準化：\par
\noindent\hspace*{1.5em} 正性標準化 \hfill IV-3\par
\noindent\hspace*{1.5em} 負性標準化 \hfill IV-3\par
\noindent\hspace*{1.5em} 二次式的標準化 \hfill IV-3，IV-4 甲\par
\noindent\hspace*{1.5em} 基底的正性標準化 \hfill IV-3\par
\noindent\hspace*{1.5em} 基底的負性標準化 \hfill IV-3\par
\noindent 鞍點 \hfill IV-4\par
\noindent 實部（複數的） \hfill V-1\par

\vspace{0.8em}
\begin{center}
\textbf{\large 十六劃}
\end{center}

\noindent 橢圓：\par
\noindent\hspace*{1.5em} 點橢圓 \hfill IV-1\par
\noindent\hspace*{1.5em} 虛橢圓 \hfill IV-1\par
\noindent 橢圓的：\par
\noindent\hspace*{1.5em} 橢圓的定義 \hfill III-1 甲\par
\noindent\hspace*{1.5em} 橢圓的標準式 \hfill II-1 乙，III-1 甲\par
\noindent\hspace*{1.5em} 橢圓的離心率 \hfill III-1 甲\par
\noindent\hspace*{1.5em} 橢圓的頂點 \hfill III-1 甲\par
\noindent\hspace*{1.5em} 橢圓的準線 \hfill II-1 乙，丙\par
\noindent\hspace*{1.5em} 橢圓的機械作圖 \hfill III-1 甲，乙\par
\noindent\hspace*{1.5em} 橢圓的特徵性質 \hfill II-1 乙\par
\noindent\hspace*{1.5em} 橢圓的切線 \hfill III-4 甲\par
\noindent\hspace*{1.5em} 橢圓的長軸 \hfill III-1 甲\par
\noindent\hspace*{1.5em} 橢圓的短軸 \hfill III-1 甲\par
\noindent\hspace*{1.5em} 橢圓的相似 \hfill III-1 乙\par
\noindent\hspace*{1.5em} 橢圓的焦點 \hfill III-1 甲\par
\noindent\hspace*{1.5em} 橢圓的中心 \hfill III-1 甲\par
\noindent 橢球面鏡的光學性質 \hfill III-4 甲\par
\noindent 錐線：\par
\noindent\hspace*{1.5em} 負類錐線（$T^-$） \hfill II-4 甲\par
\noindent\hspace*{1.5em} 零類錐線（$T^0$） \hfill II-4 乙\par
\noindent\hspace*{1.5em} 正類錐線（$T^+$） \hfill II-4 丙\par
\noindent\hspace*{1.5em} 有心錐線 \hfill IV-5\par
\noindent 歐拉（Euler）公式 \hfill V-3 甲註\par
\noindent 積化和公式 \hfill V-3 乙\par
\noindent 機械作圖：\par
\noindent\hspace*{1.5em} 橢圓的機械作圖 \hfill III-1 甲，乙\par
\noindent\hspace*{1.5em} 雙曲線的機械作圖 \hfill III-2 甲\par
\noindent\hspace*{1.5em} 拋物線的機械作圖 \hfill III-3 甲\par
\noindent 蓋射 \hfill I-5 附註\par
\noindent 點橢圓 \hfill IV-1\par

\vspace{0.8em}
\begin{center}
\textbf{\large 十七劃}
\end{center}

\noindent 壓縮性：\par
\noindent\hspace*{1.5em} 長度的壓縮性 \hfill I-3 戊\par
\noindent\hspace*{1.5em} 空間的壓縮性 \hfill I-3 戊\par
\noindent 點橢圓 \hfill IV-1\par

\vspace{0.8em}
\begin{center}
\textbf{\large 十八劃}
\end{center}

\noindent 轉化 \hfill I-3 甲\par
\noindent 轉動：\par
\noindent\hspace*{1.5em} 球的轉動 \hfill V-4 丙\par
\noindent 雙曲線：\par
\noindent\hspace*{1.5em} 共軛雙曲線 \hfill III-2 乙\par
\noindent\hspace*{1.5em} 等軸雙曲線 \hfill III-2 乙\par
\noindent 雙曲線的：\par
\noindent\hspace*{1.5em} 雙曲線的定義 \hfill III-2 甲\par
\noindent\hspace*{1.5em} 雙曲線的標準式 \hfill III-2 甲\par
\noindent\hspace*{1.5em} 雙曲線的離心率 \hfill III-2 丙\par
\noindent\hspace*{1.5em} 雙曲線的頂點 \hfill III-2 甲\par
\noindent\hspace*{1.5em} 雙曲線的準線 \hfill III-2 乙\par
\noindent\hspace*{1.5em} 雙曲線的機械作圖法 \hfill III-2 甲\par
\noindent\hspace*{1.5em} 雙曲線的特徵性質 \hfill II-1 乙\par
\noindent\hspace*{1.5em} 雙曲線的切線 \hfill III-4 乙\par
\noindent\hspace*{1.5em} 雙曲線的中心 \hfill III-2 甲\par
\noindent\hspace*{1.5em} 雙曲線的貫軸 \hfill III-2 甲\par
\noindent\hspace*{1.5em} 雙曲線的共軛軸 \hfill III-2 乙\par
\noindent\hspace*{1.5em} 雙曲線的漸近線 \hfill III-2 乙\par
\noindent\hspace*{1.5em} 雙曲線的大小 \hfill III-2 丙\par
\noindent\hspace*{1.5em} 雙曲線的開口程度 \hfill III-2 丙\par
\noindent\hspace*{1.5em} 雙曲線的焦點 \hfill III-2 甲\par
\noindent 雙曲面： \hfill II-4 甲\par
\noindent\hspace*{1.5em} 雙曲面鏡的光學性質 \hfill II-4 乙\par
\noindent 離心率： \hfill II-4\par
\noindent\hspace*{1.5em} 正類錐線的離心率 \hfill II-4 丙\par
\noindent\hspace*{1.5em} 零類錐線的離心率 \hfill II-4 乙\par
\noindent\hspace*{1.5em} 負類錐線的離心率 \hfill II-4 甲\par
\noindent\hspace*{1.5em} 橢圓的離心率 \hfill III-1 甲\par
\noindent\hspace*{1.5em} 雙曲線的離心率 \hfill III-2 丙\par
\noindent\hspace*{1.5em} 拋物線的離心率 \hfill III-3 乙\par

\vspace{0.8em}
\begin{center}
\textbf{\large 十九劃}
\end{center}

\noindent 鏡射： \hfill I-4 丙\par
\noindent\hspace*{1.5em} 鏡射軸 \hfill I-4 丙\par
\noindent\hspace*{1.5em} 以直線為對稱軸的鏡射 \hfill I-4 丙\par
\noindent\hspace*{1.5em} 對於圓的鏡射 \hfill III-6，V-2 戊\par
\noindent 鏡像點 \hfill III-6\par
\noindent 羅倫茲（Lorentz）變換 \hfill I-3 戊\par

\end{multicols}
